\documentclass[12pt]{report}

\usepackage[utf8]{inputenc}
\usepackage[french]{babel}
\usepackage{fullpage}
\usepackage{graphicx}
\usepackage{fancyhdr}	% headers/footers
\usepackage{xcolor}		% to use our own color
\usepackage{lastpage}	% to easily know the total number of pages
\usepackage{titling}	% to easily know the total number of pages
\usepackage{colortbl}	% to put color in a table background
\usepackage{datetime}	% to allow us set a new date formatting
\usepackage{multirow}   % to allow multirows in tables

% Custom defines zone

% Define useful hand-made colors
\definecolor{epiBlue}{RGB}{23,54,93}
\definecolor{lightGray}{gray}{0.85}

% Bit of code to bold an entire table row
% http://tex.stackexchange.com/questions/4811/make-first-row-of-table-all-bold
\newcolumntype{$}{>{\global\let\currentrowstyle\relax}}
\newcolumntype{^}{>{\currentrowstyle}}
\newcommand{\rowstyle}[1]{\gdef\currentrowstyle{#1}%
  #1\ignorespaces
}

% Defining a "dd/mm/yyyy" date format
\newdateformat{dashDate}{\twodigit{\THEDAY}/\twodigit{\THEMONTH}/\twodigit{\THEYEAR}}

% Define Document Title
\newcommand{\DocTitle}{[2015][Onitu][CDC]}

% end of Defines


% fancyhdr-specific commands
\setlength{\headheight}{15.2pt}

%% Defining headers and footers contents.

% Big dirty hack of the "empty" pagestyle to show header and footer on the title page (in wait of a better solution)
\fancypagestyle{empty}
{
	\renewcommand{\headrulewidth}{0pt}
	\renewcommand{\footrulewidth}{1pt}
	\fancyhead[L]{\includegraphics[height=42pt]{logo_eip.png}
	}
	\fancyhead[R]{\colorbox{epiBlue}{\color{white}\textbf{\Large{\DocTitle} } }
	}

	\fancyfoot[L]{
   		\textcolor{gray}{\{EPITECH.\}}
	}
	\fancyfoot[C]{
		\jobname
	}
	\fancyfoot[R]{}
}

\fancypagestyle{EIP}
{
	\renewcommand{\headrulewidth}{0pt}
	\renewcommand{\footrulewidth}{1pt}
	\fancyhead[L]{\includegraphics[height=42pt]{logo_eip.png}
	}
	\fancyhead[R]{\colorbox{epiBlue}{\color{white}\textbf{\Large{\DocTitle} } }
	}

	\fancyfoot[L]{
   		\textcolor{gray}{\{EPITECH.\}}
	}
	\fancyfoot[C]{
		\jobname
	}
	\fancyfoot[R]{
		\thepage/\pageref{LastPage}
	}
}

\pagestyle{EIP} % does not seem to work ...

% end of fancyhdr stuff

%Gummi|063|=)

%\title{The Title\\\normalsize A Sub-title}
\title{
	\huge{\textbf{\textcolor{epiBlue}{EIP Onitu} } }\\
	\Large{\textbf{\emph{\textcolor{gray}{Cahier des charges (CDC)} } } }
}


\begin{document}
\addtocontents{toc}{\protect\refstepcounter{page}} % makes the table of contents count pages from 1 (one)
\maketitle

\thispagestyle{empty}
\vspace*{10mm}

\textbf{\emph{\textcolor{onitu}{\large{Résumé du document} } } }\\

Résumé !!!

\clearpage


\thispagestyle{empty}
\vspace*{10mm}
\textbf{\emph{\textcolor{epiBlue}{Description du document} } }\\

\begin{tabular}{|>{\columncolor[gray]{0.85}\color{epiBlue} \bfseries } l|l|}
\hline
	Titre & \DocTitle\\
\hline
	Date & \dashDate\today \\
\hline
	Auteur & Alexandre BARON\\
\hline
	Responsable & Louis Roché\\
\hline
	E-Mail & onitu\_2015@labeip.epitech.eu\\
\hline
	Sujet & Cahier des charges\\
\hline
	Mots clés & \\
\hline
	Version du modèle & 1.0\\
\hline
\end{tabular}
\vspace*{10mm}

\textbf{\emph{\textcolor{epiBlue}{Tableau des révisions} } }\\


\begin{tabular}{|$l|p{4cm}|p{2cm}|p{5cm}|}
\hline
\rowcolor{lightGray}
\rowstyle{ \color{epiBlue} \bfseries}
	Date & Auteur & Section(s) & Commentaires\\
\hline
	03/03/2013 & Antoine Rozo & Chapitres 1 et 2 & Complétion description du sujet et description de l'environnement \\
\hline
\end{tabular}

\tableofcontents
\addtocontents{toc}{\protect\thispagestyle{empty} 
                    \protect\pagestyle{empty}}
\thispagestyle{empty}

\chapter{Rappel de l'EIP}
\thispagestyle{EIP} % seems mandatory
\setcounter{page}{1} %reset the page count

\section{Qu'est-ce qu'un EIP et Epitech}
Epitech, école d'expertise informatique en cinq ans, propose aux étudiants, à partir de leur troisième année, un projet de fin d'études: l'EIP (pour \emph{Epitech Innovative Project}).\\

À ce titre, les élèves doivent s'organiser en un groupe d'au moins cinq personnes et choisir un sujet porteur de nouveautés ou améliorant un ancien sujet. L'EIP est un passage obligatoire et unique dans la scolarité de l'étudiant, de par son envergure (18 mois) et la préparation requise. Le but est, à la fin du temps imparti, d'obtenir un projet commercialisable.


\section{Principe de base du système futur}
    Il s’agit d’un projet visant à proposer une implémentation libre et Open Source du serveur d’Ubuntu One.\\

    Ubuntu One est un service de Canonical (sponsor officiel d'Ubuntu) permettant de disposer d’un espace de stockage en ligne qui sera synchronisé entre différents ordinateurs et périphériques compatibles via un logiciel client. Le client et le protocole d’Ubuntu One sont disponibles sous licence libre. Néanmoins, le serveur est propriétaire et n’a pas été publié.\\


    Notre objectif est donc de proposer un équivalent libre de ce serveur, afin
    de profiter des fonctionnalités d’Ubuntu One par exemple à des fins d’auto-hébergement.\\

    Le client officiel n’étant pas capable d’utiliser un serveur différent de celui
    de Canonical, nous prévoyons également d'effectuer un fork de celui-ci afin d'y ajouter cette option.\\

    Des fonctionnalités d'Ubuntu One ne seront reprises que celles axées sur le stockage des fichiers en lui-même, et nous ne nous intéresserons donc pas au streaming de fichiers audio par exemple.\\

    Notre cible est l'utilisateur averti, soucieux des problématiques de centralisation des données, et son entourage, à qui il fera profiter le serveur mis en place. Il n'est pas forcément technicien mais assez curieux, le profil même de l'utilisateur Ubuntu.

\chapter{Présentation de l'environnement de réalisation}
\thispagestyle{EIP} % seems mandatory
\section{Environnement de réalisation}
Premièrement, un point important de la réalisation sera l'utilisation de Git pour l'hébergement du code, profitant d'un système de branches efficace et offrant ainsi à chacun un bon confort de développement. Ce Git étant hébergé sur github, nous bénéficions de plus de fonctionnalités annexes telles que les issues ou le wiki de façon à disposer d'un espace clair pour discuter de différents points sur projet, ou encore de rédiger de la documentation.\\

Nous n'imposerons aucun IDE, considérant que chacun est libre d'utiliser l'environnement avec lequel il sera le plus performant. Il va de paire que les fichiers de configuration liés à un tel IDE ou autres fichiers temporaires n'auront pas leur place sur le dépôt.\\

Le code sera écrit en python3, cette version n'est pas aujourd'hui la plus répandue, mais ceci est amené à changer dans les années à venir, et donc il sera communément admis lors de la finition de notre projet.\\
Ce code sera produit en respect de la PEP8, imposant diverses règles qui impliquent un code propre et agréable à la lecture, de plus très couramment usitée chez les développeurs python.\\

De plus, nous serons dans le cadre de l'EIP soumis à des contraintes de portabilité. Ces contraintes seront annihilées par python, multi-plateforme par nature, de même que l'ensemble des biliothèques et technologies que nous envisageons d'utiliser.

\section{Environnement matériel}
Notre projet consistant principalement en un serveur de stockage, il nous faudra une ou plusieurs machines faisant office de serveur pour tester un déploiement ainsi que différents postes clients afin d'en vérifier correctement et exhaustivement la validité des fonctionnalités.\\

Niveau serveur, nous pensons qu'une seule machine physique serait nécessaire, sur laquelle nous pourrons en faire tourner un certain nombre de virtuelles, trois par exemple.\\

Pour les clients, nous utiliserons principalement nos propres PC, ansi que possiblement diverses plateformes mobiles de façon à tester le plus grand étandard de clients possibles et vérifier que le serveur réagit bien en conséquences.\\

Nous pensons aussi disposer d'une machine sur laquelle faire tourner un serveur d'intégration continue dans le but de mener à bien les tests de premier niveau décrits plus bas.

\section{Architecture technique}
\section{Existing components}
\thispagestyle{EIP} % seems mandatory
Several existing clients use the Ubuntu One API, on several platforms. Onitu is compatible with those clients on the sole condition they allow the user to specify the server's address.\\

The Ubuntu One protocol, \textit{ubuntuone-storage-protocol}, uses Google's data organisation solution \textbf{Protocol Buffers}. Onitu uses the same protocol and the same tools.\\

The network comunication will mostly be using \textbf{Twisted}, a Python framework for \emph{event-driven} network programing. This framework is largely used and is in actif development.\\

Several libraries can be used to comunicate with the different services, through the \textit{drivers}. The choice of those libraries is left to the developers of the \textit{drivers}. However it must be reasonable and not cause any portability issues.

\section{Gestion de la sécurité}
\section{Points sensibles}
Un élément essentiel est le protocole d'Ubuntu One sur lequel nous nous basons. Si Canonical venait à le changer brusquement, cela constituerait un point sensible assez important. Il est peu probable que cela se produise, mais nous aurions alors le choix entre deux solutions: si les changements sont mineurs ou que notre implémentation est assez peu complète, nous pourrions nous adapter à ce nouveau protocole. Dans le cas contraire, libre oblige, nous partirions sur un fork du projet en gardant ainsi l'ancien protocole chez les clients.\\

Un autre point est la possibilité que d'autres serveurs concurrents se montent pendant le développement de notre projet.\\

On peut également se poser la question de la compatibilité multi-OS (Windows, Mac OS X, *BSD, GNU/Linux). Ce point est important car il va nous orienter dans le choix du langage utilisé et des bibliothèques ainsi que des contraintes matériels, par exemple, avoir un mac pour les tests sur Mac OS.\\

Enfin, nous serons confrontés au théorème CAP, s'appliquant aux systèmes de calcul distribué, et affirmant que trois contraintes ne peuvent simultanément être respéctées: la cohérence, la disponibilité, et la résistance au morcellement.

\chapter{Description des différentes parties du programme à réaliser}
\thispagestyle{EIP} % seems mandatory

\chapter{Description de la base de données}
\thispagestyle{EIP} % seems mandatory

\chapter{Description des tests de premier niveau}
\thispagestyle{EIP} % seems mandatory

\chapter{Organisation du projet}
\thispagestyle{EIP} % seems mandatory
Nous avons la chance de débuter le projet dans des conditions proches de ce que nous vivrons l'année prochaine, au décalage horaire près, dans le sens où nous sommes répartis sur trois sites différents. Ainsi, notre organisation prend en compte cet élément dès le départ, et nous ne subirons pas de choc, ou moindre, l'an prochain.\\

Pour cela, nous avons au début du projet mis en place un système de réunions hebdomadaires, à savoir chaque mercredi soir, où nous discutons principalement des documents à rendre ainsi que de leurs échéances. En période de crise, il nous arrive aussi de prévoir de nouvelles réunions.\\

En parallèle, nous sommes constamment présents sur un canal IRC sur lequel nous débattons des axes que nous souhaitons prendre pour notre projet. Il nous sert aussi aux préparatifs des réunions, nous y fixons en effet les heures et nous réunissons préalablement, après quoi nous passons sur le système de bulles de Google pour procéder à la réunion proprement dite.\\

De plus, notre projet étant hébergé sur Github, nous bénéficions des avantages de ce dernier dans sa gestion, à savoir principalement un \textit{wiki} sur lequel nous écrivons par exemple nos compte-rendus de réunions ou mettons à disposition des conversations importantes que nous avons pu avoir, ainsi qu'un système d'\textit{issues} et \textit{milestones}, qui nous servent à la répartition des tâches, et à avoir une vision des échéances.

\chapter{Annexes}
\thispagestyle{EIP} % seems mandatory

\end{document}
